\chapter{Using \textsc{Lark} as a Parser Generator  \label{chapter:ply}}
\epigraph{\blue{\large\calligfont Give a man a parser and he can read one programming language.} \\
          \blue{\large\calligfont Give a man a parser generator and he can read all programming languages.}}{--- Carlos Ben Yussuf}

\epigraph{\red{\large\calligfont Build a witch a fire and you keep her warm for one night.} \\
          \red{\large\calligfont Set a witch on fire and you keep her warm for the rest of her life.}}{--- Pope
        Hugo The Terrible}
        

\begin{figure}[h] 
\centering
  \includegraphics[width=17.0cm]{Abbildungen/lark.png}
\caption{The Lark parser generator in action.}
\label{fig:lark.png}
\end{figure}


Many modern programming languages, e.g.~\texttt{C}, \textsl{Java}, or \textsc{Go} can
be parsed using a so-called \simtextsc{Lalr}(1)-Parser\footnote{
  LALR(1) is an acronym that stands for \blue{Look-Ahead Left-to-Right} parsing with 1 token of lookahead.
  We will discuss the theory of \simtextsc{Lalr}(1) parsing at the end of Chapter \ref{chapter:bottom-up}.}.
Therefore, this chapter introduces the parser generator
\href{https://lark-parser.readthedocs.io/}{\simtextsc{Lark}}\index{Lark@\simtextsc{Lark}}, which can be used to
generate a parser for any language that has an \simtextsc{Lalr}(1) grammar.\footnote{
  \simtextsc{Lark} is a pure \textsl{Python} library that can be installed via the command
  \\[0.1cm]
  \hspace*{1.3cm}
  \textcolor{amethyst}{\texttt{pip install lark}}.
  \\[0.1cm]
  Besides \simtextsc{Lalr}(1) parsing, \simtextsc{Lark} also implements the \blue{Earley} algorithm, which is
  able to parse \emph{every} context-free grammar.  We will discuss the Earley algorithm in a later chapter.}
In Chapter \ref{chapter:ply-lex} we have already seen how a scanner can be implemented with the help of the
module \texttt{re}.  \simtextsc{Lark} goes one step further:  Both the \blue{lexical structure} and the
\blue{syntactical structure} of a language are described in a \emph{single} grammar, and \simtextsc{Lark}
generates the scanner \emph{and} the parser from this description.  Hence, this chapter no longer has to
distinguish between the specification of a scanner and the specification of a parser.
We will discuss two examples that use \simtextsc{Lark} as a parser generator.
\begin{enumerate}
\item First, we implement a symbolic calculator.
\item In our second example, we develop a program for symbolic differentiation.
      This example shows how a parser can be used to construct an abstract syntax tree.
\end{enumerate}
In Chapter \ref{chapter:interpreter} we will then use \simtextsc{Lark} to create an interpreter for a language
that is a small fragment of the programming language \texttt{C}.

\remarkEng
Traditional parser generators for \textsl{Python}, for example
\href{https://www.dabeaz.com/ply/}{\simtextsc{Ply}}, encode every grammar rule as a \textsl{Python}
function whose \blue{document string} contains the rule, and the \blue{semantic action}\index{semantic action} associated with a
rule is the body of this function.  Therefore, syntax and semantics are interleaved.  \simtextsc{Lark}
takes a different approach:  The grammar is written down as an ordinary string using a dialect of the
\blue{extended Backus-Naur form} \simtextsc{Ebnf}, while the semantics is implemented separately.  This
separation has two advantages:  The grammar remains readable as a grammar, and the same grammar can be
reused for different purposes.  \eox


\section{A Symbolic Calculator  \label{section:calculator}}
Figure \ref{fig:calculator.g} on page \pageref{fig:calculator.g} shows the grammar of a simple 
\blue{symbolic calculator}\index{calculator, symbolic}, i.e.~a calculator that supports the use of variables.  

\begin{figure}[!ht]

  \begin{center}    
  \framebox{
  \framebox{
  \begin{minipage}[t]{9cm}
  \begin{eqnarray*}
  \textsl{stmnt}   & \rightarrow & \;\simtextsc{Identifier} \quoted{:=} \textsl{expr} \quoted{;}\\
                   & \mid        & \;\textsl{expr} \quoted{;}                   \\[0.2cm]
  \textsl{expr}    & \rightarrow & \;\textsl{expr} \quoted{+} \textsl{product}  \\
                   & \mid        & \;\textsl{expr} \quoted{-} \textsl{product}  \\
                   & \mid        & \;\textsl{product}                           \\[0.2cm]
  \textsl{product} & \rightarrow & \;\textsl{product} \quoted{*} \textsl{factor}\\
                   & \mid        & \;\textsl{product} \quoted{/} \textsl{factor}\\
                   & \mid        & \;\textsl{factor}                            \\[0.2cm]
  \textsl{factor}  & \rightarrow &   \quoted{(} \textsl{expr} \quoted{)}        \\
                   & \mid        & \;\simtextsc{Number}                            \\
                   & \mid        & \;\simtextsc{Identifier}                        
  \end{eqnarray*}
  \vspace*{-0.5cm}
  \end{minipage}}}
  \end{center}
  \caption{A grammar for a symbolic calculator.}
  \label{fig:calculator.g}
\end{figure}

In order to generate a symbolic calculator that is based on this grammar, we have to transcribe the
grammar shown in Figure \ref{fig:calculator.g} into the \simtextsc{Ebnf} dialect that is understood by
\simtextsc{Lark}.  Figure \ref{fig:01-Calculator.ipynb:grammar} on
page \pageref{fig:01-Calculator.ipynb:grammar} shows the result.  Note that there is no separate scanner
specification:  The terminals are declared inside the very same string that also declares the grammar
rules.

\begin{figure}[!ht]
\centering
\begin{minted}[ frame        = lines, 
                framesep     = 0.3cm, 
                bgcolor      = sepia,
                numbers      = left,
                numbersep    = -0.2cm,
                xleftmargin  = 0.8cm,
                xrightmargin = 0.8cm,
              ]{python3}
    from lark import Lark, Token

    calc_grammar = """
    ?start: stmnt

    ?stmnt: IDENTIFIER ":=" expr ";" -> assign
          | expr ";"                 -> expr_stmnt

    ?expr: expr "+" product          -> add
         | expr "-" product          -> sub
         | product

    ?product: product "*" factor     -> mul
            | product "/" factor     -> div
            | factor

    ?factor: "(" expr ")"
           | NUMBER                  -> number
           | IDENTIFIER              -> var

    // Lexical definitions
    %import common.CNAME -> IDENTIFIER
    %import common.NUMBER
    %import common.WS

    ASSIGN: ":="
    SEMI:   ";"

    // Ignore whitespace
    %ignore WS
    """

    parser = Lark(calc_grammar, parser='lalr')
\end{minted}
\vspace*{-0.3cm}
\caption{The grammar of the symbolic calculator, written for \simtextsc{Lark}.}
\label{fig:01-Calculator.ipynb:grammar}
\end{figure}

\noindent
Figure \ref{fig:01-Calculator.ipynb:grammar} is discussed line by line.
\begin{enumerate}
\item Line 1 imports the class \texttt{Lark}, which is responsible for computing the parse table, and the
      class \texttt{Token}.  A \texttt{Token} represents a leaf of the parse tree, i.e.~an object that has
      been produced by the scanner.
\item Line 3 -- 31 define the string \texttt{calc\_grammar}.  Note that this is an ordinary
      \textsl{Python} string:  As far as \textsl{Python} is concerned, it has no structure at all.  Only
      when this string is handed to \simtextsc{Lark} in line 33 is it interpreted as a grammar.

      \simtextsc{Lark} distinguishes syntactical variables from tokens by their spelling:  The names of
      \blue{syntactical variables} have to be written in \textsl{lower case}, while the names of
      \blue{tokens} have to be written in \textsl{upper case}.  Therefore, \texttt{expr} is a syntactical
      variable, while \texttt{NUMBER} is a token.
\item Line 4 specifies that the syntactical variable \texttt{stmnt} is the \blue{start symbol} of the
      grammar.

      The prefix ``\texttt{?}'' that is attached to the variable \texttt{start} is an
      \blue{inlining directive}\index{inlining directive (Lark)@inlining directive (\simtextsc{Lark})}:  If a rule that is marked with ``\texttt{?}'' produces a node that has
      exactly one child, then this node is discarded and its child is used in its place.  Without this
      directive, the parse tree would contain a lot of nodes that have a single child and that carry no
      information.
\item Line 6 -- 7 define the syntactical variable \texttt{stmnt}.  This definition corresponds to the two
      grammar rules
      \\[0.2cm]
      \hspace*{1.3cm}
      $\textsl{stmnt} \rightarrow \simtextsc{Identifier} \quoted{:=} \textsl{expr} \quoted{;}$
      \quad and \quad
      $\textsl{stmnt} \rightarrow \textsl{expr} \quoted{;}$
      \\[0.2cm]
      of the grammar shown in Figure \ref{fig:calculator.g}.  Observe that the different alternatives of a
      rule are separated by the symbol ``\texttt{|}'', exactly as in the mathematical notation that we
      have used so far, and that literal strings like \squoted{:=} are simply written down as
      \textsl{Python} strings.

      The construct ``\texttt{-> assign}'' is an \blue{alias}\index{alias (Lark)@alias (\simtextsc{Lark})}.  It specifies the name of the node that is
      created in the parse tree when this alternative is used.  Without an alias, all nodes created by the
      variable \texttt{stmnt} would be called \texttt{stmnt} and we would have no easy way to distinguish
      an assignment from an expression statement.  Note further that an alias also suppresses the inlining
      that is requested by the prefix ``\texttt{?}''.  Therefore, a node of the form \texttt{expr\_stmnt}
      is created even though it has just one child.
\item Line 9 -- 15 define the variables \texttt{expr} and \texttt{product}.  These rules are literal
      transcriptions of the corresponding rules of Figure \ref{fig:calculator.g}.  In particular, they
      encode the \blue{precedence}\index{precedence!of arithmetic operators} of the arithmetic operators, and they are \blue{left recursive}.
      Left recursion is not a problem here, since \simtextsc{Lalr}(1) parsing is a \blue{bottom-up}
      technique.
\item Line 17 -- 19 define the variable \texttt{factor}.  The first alternative in line 17 does not have
      an alias.  The reason is that the tokens \squoted{(} and \squoted{)} are discarded by
      \simtextsc{Lark} and hence this alternative produces a node with a single child.  As the variable
      \texttt{factor} is marked with the prefix ``\texttt{?}'', this node is inlined and the parentheses
      leave no trace in the parse tree.  This is exactly what we want, because the only purpose of the
      parentheses is to override the precedence of the operators.
\item Line 22 -- 24 import predefined token definitions from the library \texttt{common} that is shipped
      with \simtextsc{Lark}.  Instead of writing our own regular expressions we reuse the following
      definitions:
      \begin{enumerate}
      \item \texttt{CNAME} matches an identifier in the sense of the language \texttt{C}.  The construct
            ``\texttt{-> IDENTIFIER}'' renames this token to \texttt{IDENTIFIER}.
      \item \texttt{NUMBER} matches both integers and floating point numbers.
      \item \texttt{WS} matches \blue{white space}, i.e.~blanks, tabulators, and newlines.
      \end{enumerate}
\item Line 26 -- 27 give explicit names to the tokens \squoted{:=} and \squoted{;}.  These two
      declarations do not change the language that is accepted:  Both of these tokens are still discarded
      when the parse tree is built.  Their only purpose is to ensure that these tokens have readable names
      instead of the names that \simtextsc{Lark} would generate automatically.  This is convenient once we
      start to inspect the \simtextsc{Lalr}(1) parse table.
\item Line 30 instructs the scanner to discard all white space.
\item Line 33 creates the parser.  The keyword argument \texttt{parser='lalr'} selects the
      \simtextsc{Lalr}(1) algorithm.  It is at this point that the parse table is computed.  If the given
      grammar is not an \simtextsc{Lalr}(1) grammar, then this line raises an exception.
\end{enumerate}
\FloatBarrier

\noindent
Once the parser has been created, a string \texttt{s} can be parsed by evaluating the expression
\\[0.2cm]
\hspace*{1.3cm}
\texttt{parser.parse(s)}.
\\[0.2cm]
This expression returns an object of class \texttt{lark.Tree}.  A \texttt{Tree} has two attributes:
\begin{enumerate}
\item \texttt{data} is the name of the node.  This is either the name of the syntactical variable that has
      produced this node, or, if an alias has been specified, the name given by this alias.
\item \texttt{children} is the list of the subtrees of this node.  The elements of this list are either
      objects of class \texttt{Tree} or objects of class \texttt{Token}.  As the class \texttt{Token} is a
      subclass of \texttt{str}, a token can be used like a string.
\end{enumerate}
Note that the parse tree returned by \texttt{parser.parse} does not know anything about the
\emph{meaning} of the program that has been parsed.  Providing this meaning is our job and we do it in two
steps:
\begin{enumerate}
\item First, we convert the object of class \texttt{lark.Tree} into an \blue{abstract syntax tree} that is
      represented as a \blue{nested tuple}\index{nested tuple}.
\item Second, we evaluate this nested tuple.
\end{enumerate}
Admittedly, \simtextsc{Lark} offers the class \texttt{Transformer}, which could be used to evaluate the
parse tree directly.  However, the two-step approach has the advantage that the abstract syntax tree is a
\blue{pure} \textsl{Python} object:  It does not depend on \simtextsc{Lark} in any way and we can inspect
it, print it, or store it without having to know anything about the internals of the parser.

\begin{figure}[!ht]
\centering
\begin{minted}[ frame        = lines, 
                framesep     = 0.3cm, 
                bgcolor      = sepia,
                numbers      = left,
                numbersep    = -0.2cm,
                xleftmargin  = 0.0cm,
                xrightmargin = 0.0cm,
              ]{python3}
    def ast_to_tuples(node):
        if isinstance(node, Token):
            return str(node)
        operator_map = {
            'assign':     ':=',
            'expr_stmnt': 'eval',
            'add':        '+',
            'sub':        '-',
            'mul':        '*',
            'div':        '/',
            'number':     'num',
            'var':        'var'
        }
        op       = operator_map[node.data]
        children = [ast_to_tuples(child) for child in node.children]
        return (op, *children)
\end{minted}
\vspace*{-0.3cm}
\caption{Converting a parse tree into an abstract syntax tree.}
\label{fig:01-Calculator.ipynb:ast}
\end{figure}

\noindent
Figure \ref{fig:01-Calculator.ipynb:ast} on page \pageref{fig:01-Calculator.ipynb:ast} shows the function
\texttt{ast\_to\_tuples} that performs the first of these steps.  This function is discussed next.
\begin{enumerate}
\item Line 2 -- 3 deal with the base case:  If the given node is a \texttt{Token}, then it is a leaf of the
      parse tree.  In this case we return the string that is stored in this token.
\item Line 4 -- 13 define the dictionary \texttt{operator\_map}.  This dictionary translates the node names
      that are used by \simtextsc{Lark} into the operator symbols that we want to use in our abstract
      syntax tree.  For example, the node name \texttt{add}, which has been introduced as an alias in line
      9 of Figure \ref{fig:01-Calculator.ipynb:grammar}, is translated into the string \squoted{+}.
\item Line 14 looks up the name of the given node.  
\item Line 15 converts all children of the given node recursively.
\item Line 16 returns the resulting nested tuple:  The first component is the operator, the remaining
      components are the converted children.
\end{enumerate}

\exampleEng
The string
\\[0.2cm]
\hspace*{1.3cm}
\texttt{"x := 2 * (y + 1);"}
\\[0.2cm]
is converted into the nested tuple
\\[0.2cm]
\hspace*{0.3cm}
\texttt{(':=', 'x', ('*', ('num', '2'), ('+', ('var', 'y'), ('num', '1'))))}.
\\[0.2cm]
Note the following details:
\begin{enumerate}
\item The parentheses of the input string have disappeared, since they are already reflected by the
      structure of the tuple.
\item The identifier \texttt{x} on the left hand side of the assignment operator is represented by the
      plain string \squoted{x}, while the identifier \texttt{y} is represented by the tuple
      \texttt{('var', 'y')}.  The reason is that \texttt{x} is the \emph{name} of the variable that is
      assigned to, while \texttt{y} denotes the \emph{value} of a variable and hence has to be looked up.
\item The numbers are still represented as strings.  They are converted into floating point numbers only
      when the abstract syntax tree is evaluated. \eox
\end{enumerate}

\begin{figure}[!ht]
\centering
\begin{minted}[ frame        = lines, 
                framesep     = 0.3cm, 
                bgcolor      = sepia,
                numbers      = left,
                numbersep    = -0.2cm,
                xleftmargin  = 0.0cm,
                xrightmargin = 0.0cm,
              ]{python3}
   Names2Values = {}

   def evaluate_adt(adt):
       match adt:
           case ('num', value):
               return float(value)
           case ('var', name):
               if name not in Names2Values:
                   print(f"Error: variable '{name}' is undefined.")
                   return 0.0
               return Names2Values[var_name]
           case (':=', var_name, expr):
               Names2Values[var_name] = evaluate_adt(expr)
           case ('eval', expr):
               print(evaluate_adt(expr))
           case (('+' | '-' | '*' | '/') as op, left, right):
               val_left, val_right = map(evaluate_adt, [left, right])
               match op:               
                   case '+': return val_left + val_right
                   case '-': return val_left - val_right
                   case '*': return val_left * val_right
                   case '/': return val_left / val_right
           case _:
               raise ValueError(f'Unknown ADT structure: {adt}')
\end{minted}
\vspace*{-0.3cm}
\caption{Evaluating the abstract syntax tree.}
\label{fig:01-Calculator.ipynb:eval}
\end{figure}

\noindent
Figure \ref{fig:01-Calculator.ipynb:eval} on page \pageref{fig:01-Calculator.ipynb:eval} shows the
function \texttt{evaluate\_adt} that implements the semantics of our calculator.
\begin{enumerate}
\item Line 1 initializes the dictionary \texttt{Names2Values}.  For every identifier $x$ that is defined
      interactively, \texttt{Names2Values[$x$]} is the value associated with $x$.
\item Line 4 starts a \texttt{match} statement.  \blue{Structural pattern matching}\index{pattern matching} is the natural way to
      work with an abstract syntax tree that is represented as a nested tuple:  Every case of the
      \texttt{match} statement corresponds to one of the operators that occur in our abstract syntax tree.
\item Line 5 -- 6 evaluate a number.  This is the place where the string produced by the scanner is
      converted into a floating point number.
\item Line 7 -- 11 evaluate a variable.  If the variable has not been defined previously, an error message
      is printed and the value $0.0$ is returned.
\item Line 12 -- 13 implement an assignment:  The expression on the right hand side is evaluated and the
      resulting value is stored in the dictionary \texttt{Names2Values}.  Note that an assignment does not
      return a value.
\item Line 14 -- 15 implement an expression statement:  The expression is evaluated and the resulting value
      is printed.
\item Line 16 -- 22 evaluate the arithmetic operators.  The pattern in line 16 uses an \blue{or pattern}\index{pattern matching!or pattern} in
      order to capture all four binary operators at once.  The keyword \texttt{as} binds the operator that
      has actually been found to the variable \texttt{op}.  Then, both operands are evaluated recursively
      and the inner \texttt{match} statement in line 19 -- 23 selects the operation that has to be
      performed.
\item Line 23 -- 24 catch the case that the abstract syntax tree has a shape that is not supported.  If
      this happens, we have made a mistake, either in the grammar or in the function
      \texttt{ast\_to\_tuples}.
\end{enumerate}
\FloatBarrier

\begin{figure}[!ht]
\centering
\begin{minted}[ frame        = lines, 
                framesep     = 0.3cm, 
                bgcolor      = sepia,
                numbers      = left,
                numbersep    = -0.2cm,
                xleftmargin  = 0.8cm,
                xrightmargin = 0.8cm,
              ]{python3}
    def main():
        while True:
            try:
                s = input('calc> ')
                if s == '':
                    break
                raw_tree   = parser.parse(s)
                adt_tuples = ast_to_tuples(raw_tree)
                evaluate_adt(adt_tuples)
            except Exception as e:
                print(f'Error: {e}')
\end{minted}
\vspace*{-0.3cm}
\caption{A read-eval-print loop for the symbolic calculator.}
\label{fig:01-Calculator.ipynb:main}
\end{figure}

\noindent
Finally, Figure \ref{fig:01-Calculator.ipynb:main} on page \pageref{fig:01-Calculator.ipynb:main} shows
the function \texttt{main}, which is a \blue{read-eval-print loop}\index{read-eval-print loop}:  It reads a string \texttt{s} from the
command line and then performs the three steps that we have discussed:  Line 7 parses the string, line 8
converts the resulting parse tree into an abstract syntax tree, and line 9 evaluates this tree.  The loop
is terminated as soon as the user enters an empty line.

Observe how \simtextsc{Lark} reports errors:  Instead of calling an error function that has been provided
by the user, \simtextsc{Lark} raises an \blue{exception}.  Both lexical errors and syntax errors are
signalled by subclasses of the class \texttt{lark.exceptions.UnexpectedInput}.  As these exceptions carry
the position where the error has occurred, the messages produced by \simtextsc{Lark} are usually much more
informative than the messages that we could produce ourselves with a modest amount of effort.  The
\texttt{try}-\texttt{except} statement in line 3 and line 10 -- 11 catches these exceptions so that a
single erroneous input does not terminate the calculator.

The Jupyter notebook \texttt{01-Calculator}, which is available at
\\[0.2cm]
\hspace*{1.3cm}
\href{https://github.com/karlstroetmann/Formal-Languages/blob/master/Python/Chapter-07/01-Calculator.ipynb}{Python/Chapter-07/01-Calculator.ipynb}
\\[0.2cm]
implements the symbolic calculator discussed in this section.

\remarkEng
The notebook \texttt{01-Calculator} also demonstrates how the \simtextsc{Lalr}(1) states that have been
computed by \simtextsc{Lark} can be printed.  At this point, these states will not make much sense.
However, it is worthwhile to return to this part of the notebook once we have discussed the theory of
\simtextsc{Lalr}(1) parsing in Chapter \ref{chapter:bottom-up}.  \eox

\exerciseEng
Extend the grammar shown in Figure \ref{fig:01-Calculator.ipynb:grammar} so that the calculator supports
the \blue{exponentiation operator} \squoted{**} and the \blue{unary minus}\index{unary minus}.  The exponentiation operator
binds stronger than multiplication and it is \blue{right associative}\index{right associative}, i.e.~the expression
\texttt{2 ** 3 ** 2} has the value $2^{(3^2)} = 512$.  Furthermore, extend the functions
\texttt{ast\_to\_tuples} and \texttt{evaluate\_adt} accordingly.
\\[0.2cm]
\hint
Introduce a new syntactical variable \texttt{power} between \texttt{product} and \texttt{factor} and make
this variable right recursive.

\section{Symbolic Differentiation  \label{section:differentiator}}
Our second example shows how a parser can be used to construct an \blue{abstract syntax tree} that is
afterwards transformed by a program.  Concretely, we implement \blue{symbolic differentiation}\index{symbolic differentiation}:  Given a
string \texttt{s} that represents an arithmetic expression containing the variable \texttt{x}, we want to
compute a string that represents the derivative
\\[0.2cm]
\hspace*{1.3cm}
$\ds\frac{\mathrm{d}\;}{\mathrm{d}x} \texttt{s}$.
\\[0.2cm]
For example, given the input string \texttt{"x*exp(x)"} the program should produce the output string
\\[0.2cm]
\hspace*{1.3cm}
\texttt{"exp(x) + x*exp(x)"}
\\[0.2cm]
since we have
\\[0.2cm]
\hspace*{1.3cm}
$\ds\frac{\mathrm{d}\;}{\mathrm{d}x} \bigl( x \cdot \mathrm{e}^x \bigr) = \mathrm{e}^x + x \cdot \mathrm{e}^x$.
\\[0.2cm]
In contrast to the calculator of the previous section, the abstract syntax tree is no longer evaluated
immediately.  Instead, it is the \emph{data} that our program works on:  It is transformed into a second
abstract syntax tree that represents the derivative, this tree is then simplified, and only at the very
end is it converted back into a string.  This is the reason why building an explicit abstract syntax tree
pays off.

Figure \ref{fig:differentiator.g} on page \pageref{fig:differentiator.g} shows the grammar of the
arithmetic expressions that our program supports.  Besides the four basic arithmetic operators, this
grammar supports the \blue{exponentiation operator} \squoted{**}, the \blue{unary minus}, and the two
functions $\exp$ and $\ln$.  For the sake of simplicity, \texttt{x} is the only variable.

\begin{figure}[!ht]

  \begin{center}    
  \framebox{
  \framebox{
  \begin{minipage}[t]{10cm}
  \begin{eqnarray*}
  \textsl{expr}    & \rightarrow & \;\textsl{expr} \quoted{+} \textsl{product}      \\
                   & \mid        & \;\textsl{expr} \quoted{-} \textsl{product}      \\
                   & \mid        & \;\textsl{product}                               \\[0.2cm]
  \textsl{product} & \rightarrow & \;\textsl{product} \quoted{*} \textsl{unary}     \\
                   & \mid        & \;\textsl{product} \quoted{/} \textsl{unary}     \\
                   & \mid        & \;\textsl{unary}                                 \\[0.2cm]
  \textsl{unary}   & \rightarrow &   \quoted{-} \textsl{unary}                      \\
                   & \mid        & \;\textsl{factor}                                \\[0.2cm]
  \textsl{factor}  & \rightarrow & \;\textsl{base} \quoted{**} \textsl{unary}       \\
                   & \mid        & \;\textsl{base}                                  \\[0.2cm]
  \textsl{base}    & \rightarrow &   \quoted{exp} \quoted{(} \textsl{expr} \quoted{)} \\
                   & \mid        &   \quoted{ln} \quoted{(} \textsl{expr} \quoted{)}  \\
                   & \mid        &   \quoted{(} \textsl{expr} \quoted{)}            \\
                   & \mid        & \;\simtextsc{Number}                             \\
                   & \mid        &   \quoted{x}                                     
  \end{eqnarray*}
  \vspace*{-0.5cm}
  \end{minipage}}}
  \end{center}
  \caption{A grammar for differentiable arithmetic expressions.}
  \label{fig:differentiator.g}
\end{figure}

\noindent
Figure \ref{fig:02-Differentiator.ipynb:grammar} on page \pageref{fig:02-Differentiator.ipynb:grammar}
shows the transcription of this grammar into the \simtextsc{Ebnf} dialect of \simtextsc{Lark}.

\begin{figure}[!ht]
\centering
\begin{minted}[ frame        = lines, 
                framesep     = 0.3cm, 
                bgcolor      = sepia,
                numbers      = left,
                numbersep    = -0.2cm,
                xleftmargin  = 0.8cm,
                xrightmargin = 0.8cm,
              ]{python3}
    from lark import Lark, Token

    calc_grammar = """
    ?start: expr

    ?expr: expr "+" product   -> add
         | expr "-" product   -> sub
         | product

    ?product: product "*" unary -> mul
            | product "/" unary -> div
            | unary

    ?unary: "-" unary           -> neg
          | factor

    ?factor: base "**" unary    -> pow
           | base

    ?base: "exp" "(" expr ")"    -> exp
         | "ln" "(" expr ")"     -> ln
         | "(" expr ")"
         | NUMBER                -> number
         | "x"                   -> var

    %import common.NUMBER
    %import common.WS
    %ignore WS
    """

    parser = Lark(calc_grammar, parser='lalr')
\end{minted}
\vspace*{-0.3cm}
\caption{The grammar of the differentiator, written for \simtextsc{Lark}.}
\label{fig:02-Differentiator.ipynb:grammar}
\end{figure}

\noindent
As the general structure of this grammar is the same as the structure of the grammar that we have
discussed in the previous section, we only comment on those parts that are new.
\begin{enumerate}
\item Line 4 declares the variable \texttt{expr} as the start symbol.  In contrast to the calculator,
      this language has no statements:  The input is a single expression and there is no terminating
      semicolon.
\item Line 6 -- 12 define the variables \texttt{expr} and \texttt{product} in the usual way.  Note that the
      second operand of the operators \squoted{*} and \squoted{/} is now the variable \texttt{unary}
      instead of the variable \texttt{factor}.
\item Line 14 -- 15 define the variable \texttt{unary}, which implements the \blue{unary minus}.  This rule
      is \blue{right recursive}\index{right recursive}, since a unary operator can be applied repeatedly:  The string
      \texttt{"-\,-x"} is a legal expression.
\item Line 17 -- 18 define the variable \texttt{factor}, which implements exponentiation.  The interplay
      between \texttt{unary} and \texttt{factor} deserves a closer look, because it encodes two
      conventions that every mathematician takes for granted:
      \begin{enumerate}
      \item The left operand of the operator \squoted{**} is a \texttt{base}, not a \texttt{unary} operator.
            Therefore, the string 
            \\[0.2cm]
            \hspace*{1.3cm}
            \texttt{"-x ** 2"}
            \\[0.2cm]
            is parsed as $-\bigl(x^2\bigr)$ and not as
            $(-x)^2$, i.e.~exponentiation binds stronger than the unary minus.
      \item The right operand of the operator \squoted{**} is a \texttt{unary}, not a \texttt{base}.
            Therefore, the string 
            \\[0.2cm]
            \hspace*{1.3cm}
            \texttt{"x ** -x"}
            \\[0.2cm]
            is a legal expression and, furthermore, the operator
            \squoted{**} is \blue{right associative}:  The string
            \\[0.2cm]
            \hspace*{1.3cm}
            \texttt{"x ** x ** x"}
            \\[0.2cm] is parsed as the expression $x^{\left(x^x\right)}$.
      \end{enumerate}
\item Line 20 -- 24 define the variable \texttt{base}.  The strings \texttt{"exp"} and \texttt{"ln"} in
      line 20 and line 21 are ordinary literals, i.e.~are recognized in exactly the same way as the operator
      symbols.  The nodes \texttt{exp} and \texttt{ln} have a single child.
\item Line 24 shows that the only variable of our expressions is \texttt{x}.  Since \texttt{x} is written
      as a literal string, we do not need a token for identifiers at all.
\end{enumerate}
\FloatBarrier

\noindent
Figure \ref{fig:02-Differentiator.ipynb:ast} on page \pageref{fig:02-Differentiator.ipynb:ast} shows the
function \texttt{ast\_to\_tuples} that converts the parse tree produced by \simtextsc{Lark} into a nested
tuple, together with the function \texttt{parse} that combines both steps.

\begin{figure}[!ht]
\centering
\begin{minted}[ frame        = lines, 
                framesep     = 0.3cm, 
                bgcolor      = sepia,
                numbers      = left,
                numbersep    = -0.2cm,
                xleftmargin  = 0.0cm,
                xrightmargin = 0.0cm,
                ]{python3}
    operatorTable = {'add': '+', 'sub': '-',
                     'mul': '*', 'div': '/', 'pow': '**'}
                    
    def ast_to_tuples(node):
        if isinstance(node, Token):
            return str(node)
        match node.data:
            case 'number': 
                return int(node.children[0])
            case 'var':    
                return 'x'
            case op_name:
                op_sym = operatorTable.get(op_name, op_name)
                return (op_sym, *map(ast_to_tuples, node.children))                
    
    def parse(s):
        return ast_to_tuples(parser.parse(s))
\end{minted}
\vspace*{-0.3cm}
\caption{Converting the parse tree into an abstract syntax tree.}
\label{fig:02-Differentiator.ipynb:ast}
\end{figure}


\noindent
This function serves the same purpose as the function of the same name that we have discussed in the
previous section.  However, it is organized differently:  Instead of providing a separate case for every
operator, it is \blue{table driven}.
\begin{enumerate}
\item Line 1 -- 2 define the dictionary \texttt{operatorTable}.  This dictionary maps the names of the nodes
      that are created by \simtextsc{Lark} to the operator symbols that we want to use in our abstract
      syntax tree.  For example, the node name \texttt{add}, which has been introduced as an alias in
      line 6 of Figure \ref{fig:02-Differentiator.ipynb:grammar}, is mapped to the string \squoted{+}.

      Observe that the node names \texttt{neg}, \texttt{exp}, and \texttt{ln} do \emph{not} occur in this
      dictionary.  The reason is that we want to use these names unchanged in the abstract syntax tree.
      Line 13 shows how this is achieved.
\item Line 5 -- 6 deal with the base case:  If the given node is a \texttt{Token}, then it is a leaf of
      the parse tree and hence we return the string that is stored in this token.
\item Line 7 starts a \texttt{match} statement that inspects the attribute \texttt{node.data}, i.e.~the
      name of the node.  As this name is an ordinary string, all patterns of this \texttt{match} statement
      are patterns for strings.
\item Line 8 -- 9 convert a number.  In contrast to the calculator of the previous section, the conversion
      of a string into a number happens already at this point, because the abstract syntax tree is not
      evaluated later but rather transformed symbolically.

      We will tacitly assume that all numbers are integers.
\item Line 10 -- 11 convert the variable \texttt{x}.  Since the string \texttt{"x"} occurs as an anonymous
      literal in the grammar, it is discarded when the parse tree is built.  Therefore, the node
      \texttt{var} has no children at all and we simply return the string \squoted{x}.
\item Line 12 is the most interesting line of this function.  The pattern \texttt{op\_name} is a
      \blue{capture pattern}\index{pattern matching!capture pattern}:  As \texttt{op\_name} is a plain variable name, this pattern matches
      \emph{any} value and binds this value to the variable \texttt{op\_name}.  Hence, this case is the
      default case that deals with all remaining nodes, i.e.~with all nodes that represent an operator.
\item Line 13 translates the name of the node into an operator symbol.  The second argument of the method
      \texttt{get} is the value that is returned if the given key is not present in the dictionary.  Here,
      this default value is the name of the node itself.  This is the trick that has been announced above:
      Every node name that is not listed in \texttt{operatorTable} is used unchanged.  Consequently, the
      nodes \texttt{neg}, \texttt{exp}, and \texttt{ln} are converted into tuples that start with the
      strings \squoted{neg}, \squoted{exp}, and \squoted{ln}, respectively.
\item Line 14 converts all children recursively and prepends the operator symbol.  The expression
      \\[0.2cm]
      \hspace*{1.3cm}
      \texttt{map(ast\_to\_tuples, node.children)}
      \\[0.2cm]
      applies the function \texttt{ast\_to\_tuples} to every child of the given node, while the operator
      ``\texttt{*}'' unpacks the resulting iterator into the tuple that is created.  Note that this works
      for unary and for binary operators alike:  The number of components of the resulting tuple is
      determined by the number of children of the node.
\item Line 16 -- 17 define the function \texttt{parse}, which combines both steps:  It parses the string
      \texttt{s} and immediately converts the resulting parse tree into an abstract syntax tree.
\end{enumerate}

\exampleEng
The string \texttt{"ln(x ** x) + exp(x * -x)"} is converted into the nested tuple
\\[0.2cm]
\hspace*{0.3cm}
\texttt{('+', ('ln', ('**', 'x', 'x')), ('exp', ('*', 'x', ('neg', 'x'))))}.
\hspace*{\fill} $\diamond$

\remarkEng
The table driven organization of the function \texttt{ast\_to\_tuples} pays off as soon as the language is
extended.  In order to support an additional operator or an additional function, it is usually sufficient
to extend the grammar:  If the alias of the new grammar rule is already the symbol that we want to use in
the abstract syntax tree, then the function \texttt{ast\_to\_tuples} does not have to be changed at all.
Otherwise, a single entry has to be added to the dictionary \texttt{operatorTable}.  \eox
\vspace*{0.1cm}

\noindent
Now that we can turn a string into an abstract syntax tree, we are ready to implement the function
\texttt{diff}.  This function takes an abstract syntax tree $e$ and computes an abstract syntax tree that
represents the derivative of $e$ with respect to the variable \texttt{x}.  Figure
\ref{fig:02-Differentiator.ipynb:diff} on page \pageref{fig:02-Differentiator.ipynb:diff} shows the
implementation.  Every case of the \texttt{match} statement corresponds to one of the differentiation
rules that are taught in any calculus course.

\begin{figure}[!ht]
\centering
\begin{minted}[ frame        = lines, 
                framesep     = 0.3cm, 
                bgcolor      = sepia,
                numbers      = left,
                numbersep    = 0.4cm,
                xleftmargin  = 0.0cm,
                xrightmargin = 0.0cm,
              ]{python3}
def diff(e):
    "differentiate the expression e with respect to the variable x"
    match e:
        case ('+', f, g):
            fs, gs = diff(f), diff(g)
            return ('+', fs, gs)
        case ('-', f, g):
            fs, gs = diff(f), diff(g)
            return ('-', fs, gs)
        case ('*', f, g):
            fs, gs = diff(f), diff(g)
            return ('+', ('*', fs, g), ('*', f, gs))
        case ('/', f, g):
            fs, gs = diff(f), diff(g)
            return ('/', ('-', ('*',fs,g), ('*',f,gs)), ('*',g,g))
        case ('**', f, g):
            return diff(('exp', ('*', g, ('ln', f))))
        case ('neg', f):
            return ('neg', diff(f))
        case ('ln', f):
            fs = diff(f)
            return ('/', fs, f)
        case ('exp', f):
            fs = diff(f)
            return ('*', fs, e)
        case 'x':
            return 1
    return 0
\end{minted}
\vspace*{-0.3cm}
\caption{Symbolic differentiation.}
\label{fig:02-Differentiator.ipynb:diff}
\end{figure}

\noindent
The function \texttt{diff} is discussed next.
\begin{enumerate}
\item Line 4 -- 6 implement the \blue{sum rule}\index{differentiation!sum rule}
      $\ds\frac{\mathrm{d}\;}{\mathrm{d}x}\bigl(f(x) + g(x)\bigr) 
       = \frac{\mathrm{d}\;}{\mathrm{d}x} f(x) + \frac{\mathrm{d}\;}{\mathrm{d}x} g(x)$,
      \\[0.1cm]
      while line 7 -- 9 implement the corresponding rule for the difference of two expressions.
\item Line 10 -- 12 implement the \blue{product rule}\index{differentiation!product rule}
      \\[0.2cm]
      \hspace*{1.3cm}
      $\ds\frac{\mathrm{d}\;}{\mathrm{d}x}\bigl(f(x) \cdot g(x)\bigr)
       = \Bigl(\frac{\mathrm{d}\;}{\mathrm{d}x} f(x)\Bigr) \cdot g(x)
       + f(x) \cdot \Bigl(\frac{\mathrm{d}\;}{\mathrm{d}x} g(x)\Bigr)$.
\item Line 13 -- 15 implement the \blue{quotient rule}\index{differentiation!quotient rule}
      \\[0.2cm]
      \hspace*{1.3cm}
      $\ds\frac{\mathrm{d}\;}{\mathrm{d}x}\left(\frac{f(x)}{g(x)}\right)
       = \frac{\Bigl(\frac{\mathrm{d}}{\mathrm{d}x} f(x)\Bigr) \cdot g(x)
              - f(x) \cdot \Bigl(\frac{\mathrm{d}}{\mathrm{d}x} g(x)\Bigr)}{g(x) \cdot g(x)}$.
\item Line 16 -- 17 deal with exponentiation.  Instead of implementing a separate rule for the expression
      $f(x)^{g(x)}$, we use the identity
      \\[0.2cm]
      \hspace*{1.3cm}
      $f(x)^{g(x)} = \exp\bigl(g(x) \cdot \ln(f(x))\bigr)$
      \\[0.2cm]
      and differentiate the resulting expression instead.  This is an instance of a general and very
      useful strategy:  Rather than adding a new rule, we \blue{rewrite} the expression so that it can be
      handled by the rules that we already have.  Note that this identity is only valid if $f(x)$ is
      positive.  Our implementation ignores this side condition.
\item Line 18 -- 19 implement the rule
      $\ds\frac{\mathrm{d}\;}{\mathrm{d}x}\bigl(-f(x)\bigr) 
       = -\frac{\mathrm{d}\;}{\mathrm{d}x} f(x)$.
\item Line 20 -- 22 and line 23 -- 25 implement the \blue{chain rule}\index{differentiation!chain rule} for the two functions that our
      language provides:
      \\[0.2cm]
      \hspace*{1.3cm}
      $\ds\frac{\mathrm{d}\;}{\mathrm{d}x}\ln\bigl(f(x)\bigr) 
       = \frac{\frac{\mathrm{d}}{\mathrm{d}x} f(x)}{f(x)}$
      \qquad and \qquad
      $\ds\frac{\mathrm{d}\;}{\mathrm{d}x}\exp\bigl(f(x)\bigr)
       = \Bigl(\frac{\mathrm{d}\;}{\mathrm{d}x} f(x)\Bigr) \cdot \exp\bigl(f(x)\bigr)$.
      \\[0.2cm]
      Observe that line 25 reuses the expression \texttt{e} instead of building the tree
      \texttt{('exp', f)} a second time.
\item Line 26 -- 27 implement the base case
      $\ds\frac{\mathrm{d}x}{\mathrm{d}x} = 1$.
\item Line 28 is reached if none of the previous patterns did match.  In this case, the expression does
      not contain the variable \texttt{x} and therefore it is a constant.  As the derivative of a constant
      is $0$, the value $0$ is returned.
\end{enumerate}
\FloatBarrier

\noindent
The abstract syntax tree that is computed by the function \texttt{diff} is correct, but it is unreadable.
For example, differentiating the string \texttt{"x * exp(x)"} yields a tree that corresponds to the
string
\\[0.2cm]
\hspace*{1.3cm}
\texttt{"1*exp(x) + x*1*exp(x)"}.
\\[0.2cm]
Therefore, two more functions are needed.  Since neither of them has anything to do with parsing, we only
describe what they do and refer the interested reader to the notebook for the implementation.
\begin{enumerate}
\item The function \texttt{simplify} takes an abstract syntax tree $e$ and returns an abstract syntax tree
      that is mathematically equivalent to $e$ but, in general, smaller.  The function works
      \blue{bottom-up}:  It simplifies all subexpressions first and only then tries to apply a
      simplification rule to the resulting expression.  Some of the rules produce a new expression that
      can be simplified further.  In these cases, the function \texttt{simplify} is applied recursively to
      the result.

      The rules that are implemented are listed below.  Here, $a$, $b$, and $f$ denote arbitrary
      expressions, while $c$ and $d$ denote numbers.
      \begin{enumerate}
      \item \textbf{Constant folding}:  If both operands of one of the operators
            $\texttt{+}$, $\texttt{-}$, $\texttt{*}$, $\texttt{/}$, or $\texttt{**}$ are numbers, then the
            operation is carried out.  A division is only performed if the divisor is different from $0$.
      \item \textbf{Negation}:
            $-(-a) \rightarrow a$, \quad
            $-0 \rightarrow 0$, \quad
            $-c \rightarrow d$ \; where $d$ is the negated number $c$.
      \item \textbf{Addition and subtraction}:
            \\[0.1cm]
            $a + 0 \rightarrow a$, \quad
            $0 + a \rightarrow a$, \quad
            $a - 0 \rightarrow a$, \quad
            $0 - a \rightarrow -a$, \quad
            $a - a \rightarrow 0$,
            \\[0.1cm]
            $a + (-b) \rightarrow a - b$, \quad
            $(-a) + b \rightarrow b - a$, \quad
            $a - (-b) \rightarrow a + b$.
      \item \textbf{Multiplication}:
            \\[0.1cm]
            $a \cdot 0 \rightarrow 0$, \quad
            $0 \cdot a \rightarrow 0$, \quad
            $a \cdot 1 \rightarrow a$, \quad
            $1 \cdot a \rightarrow a$, \quad
            $(-1) \cdot a \rightarrow -a$, \quad
            $a \cdot (-1) \rightarrow -a$,
            \\[0.1cm]
            $\ds a \cdot \frac{1}{a} \rightarrow 1$, \quad
            $\ds (-a) \cdot \frac{1}{a} \rightarrow -1$, \quad
            $\ds a \cdot \Bigl(-\frac{1}{a}\Bigr) \rightarrow -1$.
      \item \textbf{Division}:
            $\ds \frac{0}{a} \rightarrow 0$, \quad
            $\ds \frac{a}{1} \rightarrow a$, \quad
            $\ds \frac{a}{a} \rightarrow 1$.
      \item \textbf{Powers}:
            $a^0 \rightarrow 1$, \quad
            $a^1 \rightarrow a$, \quad
            $0^a \rightarrow 0$, \quad
            $1^a \rightarrow 1$.
      \item \textbf{Logarithm and exponential function}:
            \\[0.1cm]
            $\ln(1) \rightarrow 0$, \quad
            $\ln\bigl(\exp(a)\bigr) \rightarrow a$, \quad
            $\exp(0) \rightarrow 1$, \quad
            $\exp\bigl(\ln(a)\bigr) \rightarrow a$,
            \\[0.1cm]
            $\exp\bigl(a \cdot \ln(f)\bigr) \rightarrow f^{a}$.
            \\[0.1cm]
            The last of these rules undoes the rewriting that the function \texttt{diff} has performed in
            order to differentiate the operator \squoted{**}.
      \item \textbf{Derivatives of powers}:  These rules clean up the expressions that are produced when a
            power is differentiated.  For a number $c$ we have
            \\[0.1cm]
            $\ds a \cdot \frac{1}{f} \cdot f^{c} \rightarrow a \cdot f^{\,c-1}$
            \qquad and \qquad
            $\ds \frac{1}{f} \cdot f^{c} \rightarrow f^{\,c-1}$.
            \\[0.1cm]
            In the special case $c = 2$ the results are simplified further to $a \cdot f$ and $f$,
            respectively.
      \end{enumerate}
\item The function \texttt{toString} takes an abstract syntax tree $e$ and returns a string that
      represents $e$.  The interesting part of this function is that it does not insert parentheses
      blindly:  A subexpression is only enclosed in parentheses if the precedence of its operator is lower
      than the precedence that is required at this position.  To this end, the function uses the
      precedence table
      \\[0.2cm]
      \hspace*{1.3cm}
      \begin{tabular}{lcl}
        \squoted{+}, \squoted{-} & : & $1$, \\
        \squoted{*}, \squoted{/} & : & $2$, \\
        unary minus              & : & $3$, \\
        \squoted{**}             & : & $4$, \\
        numbers and \texttt{x}   & : & $5$.
      \end{tabular}
      \\[0.2cm]
      Since the operators \squoted{-} and \squoted{/} are not associative, the right operand of these
      operators is parenthesized already when its operator has the \emph{same} precedence.  For the
      operator \squoted{**} it is the other way round:  As this operator is right associative, its right
      operand needs no parentheses if it is a power itself, while its left operand is parenthesized
      whenever it is not an atom.  Therefore, the expression $x^{\left(x^x\right)}$ is printed as
      \texttt{x**x**x}, while $(x+1)^2$ is printed as \texttt{(x + 1)**2}.
\end{enumerate}

\noindent
Putting everything together, the derivative of a string \texttt{s} is computed by the expression
\\[0.2cm]
\hspace*{1.3cm}
\texttt{toString(simplify(diff(parse(s))))}.
\\[0.2cm]
The following table shows the results that are produced for a number of inputs.
\begin{center}
\begin{tabular}{|l|l|}
\hline
\textbf{input \texttt{s}} & \textbf{computed derivative} \\
\hline
\texttt{x*x}           & \texttt{x + x}                                        \\
\texttt{x/x}           & \texttt{0}                                            \\
\texttt{x * exp(x)}    & \texttt{exp(x) + x*exp(x)}                            \\
\texttt{-x ** 2}       & \texttt{-2*x}                                         \\
\texttt{x ** -x}       & \texttt{(-1 - ln(x))*x**(-x)}                         \\
\texttt{1/(1+exp(-x))} & \texttt{exp(-x)/((1 + exp(-x))*(1 + exp(-x)))}        \\
\texttt{ln(x ** x)}    & \texttt{(ln(x) + 1)*x**x/x**x}                        \\
\hline
\end{tabular}
\end{center}

\remarkEng
The last two rows of this table show the limits of our simplifier.  The derivative of \texttt{x*x} is
computed as \texttt{x + x} instead of \texttt{2*x}, and the derivative of \texttt{ln(x ** x)} is not
reduced to \texttt{ln(x) + 1}, although the factor $x^x/x^x$ obviously cancels.  The reason is that our
function \texttt{simplify} is nothing but a finite collection of \blue{rewrite rules}\index{rewrite rule}:  It does not bring
an expression into a \blue{normal form}\index{normal form} and therefore it cannot recognize that two syntactically different
expressions are mathematically equal.  Computing normal forms of algebraic expressions is a hard problem
that is studied in the field of \blue{computer algebra}\index{computer algebra}.  \eox

The Jupyter notebook \texttt{02-Differentiator}, which is available at
\\[0.2cm]
\hspace*{1.3cm}
\href{https://github.com/karlstroetmann/Formal-Languages/blob/master/Python/Chapter-07/02-Differentiator.ipynb}{Python/Chapter-07/02-Differentiator.ipynb}
\\[0.2cm]
implements the symbolic differentiator that has been discussed in this section.

\exerciseEng
Extend the differentiator so that it supports the trigonometric functions $\sin$ and $\cos$.  You have to
\begin{enumerate}
\item add two alternatives to the variable \texttt{base} in the grammar shown in Figure
      \ref{fig:02-Differentiator.ipynb:grammar},
\item add the corresponding cases to the function \texttt{ast\_to\_tuples}, and
\item add the rules
      \\[0.2cm]
      \hspace*{1.3cm}
      $\ds\frac{\mathrm{d}\;}{\mathrm{d}x}\sin\bigl(f(x)\bigr) 
       = \Bigl(\frac{\mathrm{d}\;}{\mathrm{d}x} f(x)\Bigr) \cdot \cos\bigl(f(x)\bigr)$
      \quad and \quad
      $\ds\frac{\mathrm{d}\;}{\mathrm{d}x}\cos\bigl(f(x)\bigr) 
       = -\Bigl(\frac{\mathrm{d}\;}{\mathrm{d}x} f(x)\Bigr) \cdot \sin\bigl(f(x)\bigr)$
      \\[0.2cm]
      to the function \texttt{diff}.
\end{enumerate}
Which changes are necessary in the function \texttt{toString}? \eox

\section{Check your Understanding}
\begin{enumerate}[(a)]
\item How does \simtextsc{Lark} distinguish the names of \blue{syntactical variables} from the names of
      \blue{tokens}?
\item What is the purpose of the prefix ``\texttt{?}'' that can be attached to a syntactical variable, and
      under which circumstances does this prefix have no effect?
\item What is an \blue{alias}, i.e.~what is the effect of the construct ``\texttt{-> name}'' that can be
      attached to an alternative of a grammar rule?
\item Which parts of a grammar rule are discarded by \simtextsc{Lark} when the parse tree is built?
\item Would you be able to write a \simtextsc{Lark} grammar for propositional formulas together with a
      function that converts the resulting parse tree into nested tuples?
\item How does \simtextsc{Lark} report lexical errors and syntax errors?
\end{enumerate}


%%% Local Variables: 
%%% mode: latex
%%% TeX-master: "formal-languages.tex"
%%% End:
